% !TeX root = ../教材题库_全部条目_转写工作稿.tex
% Source: XJ-B1-C05-RV-P208-Y0927
% 第5章 / 章末复习 复习题五 / 教材 p.208 / PDF 第215页
% OCR未校对：公式与文字均可能有误；未完成数学TeX转写；禁止直接发给学生
% 原题及相邻原解截图（核对用）：../题目原图/XJ-B1-C05-RV-P208-Y0927-01.webp
% 本文件只在转写完成后作为题干使用；OCR 条目尚不是可靠题干。
\begin{bankitem}{复习题五 第21题}{第5章 / 章末复习 复习题五 / 教材 p.208 / PDF 第215页}{XJ-B1-C05-RV-P208-Y0927}{OCR未校对：公式与文字均可能有误；未完成数学TeX转写；禁止直接发给学生}
\begin{ocrdraft}
[教材p.208 / PDF第215页]
\par
21.
\par
（多选题）已知函数丆0）一一（sinx十cos到一
\par
srnx\ocrglyph{00AF}cosxl(x\ensuremath{\text{[OCR:EUR]}}R),则
\par
2
\par
2
\par
函数不x）的最小正周期为2丌
\par
函数丆0）的值域为一1，
\par
(B)
\par
函数不x）为奇函数
\par
(C)
\par
函数丆）的图象关于直线一对称
\end{ocrdraft}
\end{bankitem}
